% Auto-generated by make_paper_assets.py
\begin{table}[t]
\centering
\caption{\label{tab:parity}%
Byte-level parity of vardictcpp against VarDict-Java 1.8.3 on five whole-exome sequencing
samples (hg19, per-sample covered-target BED, minimum allele frequency $0.01$). Sorted output
was compared line-for-line against single-threaded VarDictJava (multi-threaded Java output is
non-deterministic). \emph{Byte-identical}: rows matching the Java reference character-for-character;
\emph{FP}: rows emitted only by vardictcpp; \emph{FN}: rows emitted only by VarDictJava. All
91\,288 variant rows are byte-identical (0 FP, 0 FN).}
\begin{tabular}{lrrrrr}
\toprule
Sample & Target regions & Variant rows & Byte-identical & FP & FN \\
\midrule
\texttt{SRR15006386} & 11\,732 & 4\,014 & 4\,014\,(100\%) & 0 & 0 \\
\texttt{SRR15006375} & 12\,292 & 4\,056 & 4\,056\,(100\%) & 0 & 0 \\
\texttt{SRR15006376} & 105\,949 & 12\,739 & 12\,739\,(100\%) & 0 & 0 \\
\texttt{SRR15006540} & 135\,145 & 16\,031 & 16\,031\,(100\%) & 0 & 0 \\
\texttt{SRR8657348} & 773\,667 & 54\,448 & 54\,448\,(100\%) & 0 & 0 \\
\midrule
Total & & 91\,288 & 91\,288\,(100\%) & 0 & 0 \\
\bottomrule
\end{tabular}
\end{table}

\begin{table}[t]
\centering
\caption{\label{tab:perf}%
Runtime and peak memory of vardictcpp versus VarDict-Java 1.8.3 on three whole-exome samples
(106--774\,k target regions; hg19, minimum allele frequency $0.01$) at 1 and 8 threads, CPU-pinned.
Median of 5 replicate runs; wall clock and peak resident set size (RSS) from
\texttt{/usr/bin/time -v}, VarDictJava on JDK~25 with \texttt{-Xmx\,8g}. \emph{Speedup} $=$
Java\,$/$\,cpp runtime; $\times$\emph{less} $=$ Java\,$/$\,cpp peak RSS.
Per-run values are plotted in Figs.~\ref{fig:runtime} and \ref{fig:memory}.}
\begin{tabular}{lr rr r rr r}
\toprule
 & & \multicolumn{3}{c}{Runtime (s)} & \multicolumn{3}{c}{Peak RSS (MB)} \\
\cmidrule(lr){3-5}\cmidrule(lr){6-8}
Sample & Threads & vardictcpp & VarDictJava & Speedup & vardictcpp & VarDictJava & $\times$less \\
\midrule
\texttt{SRR15006376} (105\,949) & 1 & 44.5 & 181.1 & 4.1$\times$ & 74 & 1274 & 17$\times$ \\
 & 8 & 7.8 & 85.0 & 11.0$\times$ & 358 & 2633 & 7$\times$ \\
\addlinespace[2pt]
\texttt{SRR15006540} (135\,145) & 1 & 61.7 & 226.3 & 3.7$\times$ & 76 & 1084 & 14$\times$ \\
 & 8 & 10.4 & 96.9 & 9.3$\times$ & 312 & 2569 & 8$\times$ \\
\addlinespace[2pt]
\texttt{SRR8657348} (773\,667) & 1 & 98.0 & 473.5 & 4.8$\times$ & 122 & 977 & 8$\times$ \\
 & 8 & 14.8 & 142.4 & 9.6$\times$ & 252 & 2084 & 8$\times$ \\
\addlinespace[2pt]
\bottomrule
\end{tabular}
\end{table}

\begin{table}[t]
\centering
\caption{\label{tab:accuracy}%
Variant-calling accuracy of VarDict on the GIAB HG002 chromosome-20 exome (Agilent SureSelect
v5, Oslo University Hospital), against the NIST v4.2.1 benchmark via \texttt{rtg vcfeval} over
callable confident regions (exome depth $\geq 20$; 2.31\,Mbp; 2\,443 truth variants). Calls are
the germline operating point (PASS, allele frequency $\geq 0.2$). \textbf{vardictcpp and
VarDictJava emit byte-identical VCFs}, so every value holds for both; per-region distributions
are in Fig.~\ref{fig:accuracy}.}
\begin{tabular}{lrrrrrrr}
\toprule
Class & Truth & TP & FP & FN & Precision & Recall & F$_1$ \\
\midrule
All & 2\,443 & 2\,342 & 39 & 83 & 0.984 & 0.966 & 0.975 \\
SNV & 2\,210 & 2\,153 & 9 & 57 & 0.996 & 0.974 & 0.985 \\
Indel & 233 & 171 & 24 & 62 & 0.877 & 0.734 & 0.799 \\
\bottomrule
\end{tabular}
\end{table}
